% ---
% id: foundations/rem:computed-completeness
% kind: rem
% title: "A complete computed sum from an incomplete family"
% summary: >-
%   Three positive binary64 numbers can have computed sum exactly one while
%   their true sum exceeds one. A completeness test on the computed sum alone
%   would admit them, which is why the rounds-nowhere condition belongs to
%   exact admissibility.
% prominence: standard
% status: accepted
% uses: [def:binary64, def:exact-admissibility, lem:exact-accumulation, lem:binary64-lattice]
% assumes: []
% refs: []
% provenance: newly-derived
% lean:
% ---

\begin{remark}\label{rem:computed-completeness}
Conditions (H) and (P) of Definition~\ref{def:exact-admissibility}, together
with the identity of the \emph{computed} declared-order sum with the identity,
do not imply (C). Take the one-block map of dimension one and the family
$(1,2^{-54},2^{-54})$ in this declared order. Its entries are
\tref{binary64 numbers},
Hermitian and positive. By Lemma~\ref{lem:binary64-lattice}\,(ii) the binary64
numbers in $[1,2)$ are the multiples of $2^{-52}$, so $1+2^{-54}$ lies strictly
between the adjacent binary64 numbers $1$ and $1+2^{-52}$ and is nearer to $1$;
hence $\mathrm{fl}(1+2^{-54})=1$ (Definition~\ref{def:binary64}\,(ii)), and
the \tref{computed sequential sum} is
$0,1,1,1$. The exact sum is $1+2^{-53}\neq1$, so (C) fails and the family is
not a positive operator-valued measure. Condition (S) fails at the second
addition; it is the condition under which the computed sum is the exact sum at
every step (Lemma~\ref{lem:exact-accumulation}).
\end{remark}
